\section{Count frames before counting samples}
A sample is one channel's amplitude code at one instant. A frame groups the codes for
all channels at that instant. In a mono file those counts coincide. In a stereo file,
four frames contain eight sample codes. At sixteen bits per sample, those four stereo
frames occupy sixteen bytes. Advancing by one frame requires four bytes, whereas
advancing to the next channel inside that frame requires two.

The worked file begins with left and right codes 12000 and minus 12000. Its next frame
contains zero and 16000. In memory the order is $L_0,R_0,L_1,R_1,\ldots$, where each
subscript is a frame index. The channel labels identify different sequences; they are
not two consecutive time steps in a single sequence. Feeding that entire interleaved
array to a mono FIR would alternate between channels inside each eight-sample window.

At 16000 frames per second, a frame interval is $1/16000$ second. Four frames account
for $4/16000$ second of stored audio. Adding a second channel doubles the byte rate at
that same frame rate; it does not double the duration. Write the unit beside each
counter: frames, samples, PCM bytes and complete file bytes require different strides.
